% =====================================================================
%  Macros for Agreement Machinery
%  Every \newcommand here is read by the site and expanded exactly as
%  LaTeX would. Edit freely; you can also paste your whole preamble in.
% =====================================================================

% ---------- from the manuscript preamble ----------
\newcommand{\Next}[1][]{(\nextx#1)\xspace}
\newcommand{\Last}[1][]{(\lastx#1)\xspace}
\newcommand{\NNext}[1][]{(\anextx#1)\xspace}
\newcommand{\LLast}[1][]{(\blastx#1)\xspace}
\newcommand{\LLLast}[1][]{(\bblastx#1)\xspace}

\newcommand{\gref}[2][]{(\getref{#2}#1)\xspace}
\newcommand{\dref}[1]{\deftagex{#1}}

\newcommand{\nothing}{$\emptyset$}

\newcommand{\Lb}[1]{$\text{[}_{\text{#1}}$ }
\newcommand{\Rb}[1]{$\text{]}_{\text{#1}}$ }
\newcommand{\1}{\rlap{$'$}\xspace}
\newcommand{\0}{\rlap{\textsuperscript{$ˆ{\circ}$}}\xspace}
\newcommand{\gap}{\underline{\hspace{1.2em}}}
\newcommand{\vP}{\emph{v}P\xspace}
\newcommand{\lv}{\emph{v}\xspace}
\newcommand{\Abar}{A$'$-}
\newcommand{\pro}{\emph{pro}\xspace}
\newcommand{\subs}[1]{\textsubscript{#1}}
\newcommand{\supers}[1]{\textsuperscript{#1}}
\newcommand{\spells}{$\Longleftrightarrow$} % spellout
\newcommand{\tr}[1]{\textit{t}\textsubscript{\textit{#1}}}
\newcommand{\tht}{$\theta$}
\newcommand{\nt}{$\varnothing$}
\newcommand{\lb}{$\llbracket$}
\newcommand{\rb}{$\rrbracket$}

\newcommand{\ph}{$\phi$\xspace}
\newcommand{\gm}{$\gamma$\xspace}
\newcommand{\<}{$\langle$}
\renewcommand{\>}{$\rangle$}

% ---------- inferred from usage: check these against your preamble ----------
\newcommand{\rt}[1]{√#1}   % a root:  \rt{tpz} → √tpz
\newcommand{\sca}[1]{\textsc{#1}}         % small caps on an attached morpheme (\sca{def}-sea)
\newcommand{\xm}{\ding{55}}               % ✗ marks a blocked dependency
\newcommand{\round}[1]{(#1)}              % parenthesised label on an arrow: \round{@PF} → (@PF)

% ---------- site helpers (safe to use in LaTeX too, with the definitions below) ----------
% a feature bundle printed small under a tree node:   [T\feat{u\ph: \_\_}]
% (\textfeat is built into the site; for PDF output define it as e.g. {\scriptsize\texttt{#1}})
\newcommand{\feat}[1]{\\\textfeat{#1}}
% \headchain{ C (you) [*T *Aux Vc]{should} [V]{eat} }
%   bare head        a head with no chain line
%   [a *b c]{word}   a head chain; its lines meet under the starred head(s) (default: the first),
%                    where the chain's morphology is pronounced as {word}
%   (word)           an ordinary word on the bottom line
% (For PDF output, \headchain needs its own TikZ definition; the site draws it natively.)
